\section{Tipos de biomoléculas}

Clases de bioelementos:

\begin{itemize}
    \item Primarios: C, H, O, N, P, S
    \item Secundarios: Ca, Na, K, Mg
    \item Oligoelementos: Mn, Fe
\end{itemize}

\subsection{Inorgánicas}

Son sales minerales, y se clasifican en 2:

\begin{itemize}
    \item Solubles: Ayudan a mantener la ósmosis equilibrada, regulan el paso de solutos en la célula y regulan funciones cardiacas.
    \item Insolubles: Forman estructuras de protección, esqueletos internos, dientes, otolitos, y se acumulan como productos residuales del metabolismo.
\end{itemize}

\subsection{Orgánicas}

Las biomoléculas orgánicas son compuestos basados en carbono que forman la estructura de los seres vivos. Se clasifican en cuatro grupos principales:

\begin{itemize}
    \item \textbf{Proteínas}: Polímeros de aminoácidos
    \item \textbf{Carbohidratos}: Polímeros de monosacáridos
    \item \textbf{Lípidos}: Moléculas hidrofóbicas
    \item \textbf{Ácidos nucleicos}: Polímeros de nucleótidos
\end{itemize}

\section{Características y funciones de las biomoléculas}

    \input{tables/tab:funciones_de_cada_biomolécula}

\subsection{Importancia biotecnológica}
\begin{itemize}
    \item Proteínas recombinantes (insulina, hormonas)
    \item Enzimas industriales
    \item Bioplásticos (a partir de carbohidratos)
    \item Vacunas de ADN/ARN
\end{itemize}